\documentclass[12pt]{article}
\usepackage[utf8]{inputenc}
\usepackage[spanish]{babel}
\usepackage{amsmath}
\usepackage{amsfonts}
\usepackage{amssymb}
\usepackage{geometry}
\geometry{letterpaper, margin=2.5cm}

\begin{document}

\begin{center}
    \Large \textbf{Evaluación de Examen 1: Unidad I} \\
    \large Física para Ciencias de la Salud (005-1713) \\
    \vspace{0.5cm}
    \textbf{Estudiante:} María Meléndez \quad \textbf{C.I.:} 33.685.873
    \vspace{0.3cm} \\
    \large \textbf{Calificación Final: 8/10 puntos}
\end{center}

\vspace{0.5cm}

\section*{Análisis de Resultados}

\subsection*{1. Ángulo plano (Conversión a radianes)}
\textbf{Procedimiento y resultado:} La estudiante utiliza una regla de tres basándose en la equivalencia $90^\circ = \frac{\pi}{2}$ rad. Aunque la escritura de los pasos intermedios mezcla decimales en el numerador de una fracción $\left(\frac{37,5\pi}{90}\right)$, la lógica matemática aplicada es correcta y la lleva al resultado válido de $0,41\pi$ rad (aproximación de $\frac{5\pi}{12}$). \\
\textbf{Veredicto:} \textbf{Correcto} (2/2 pts).

\subsection*{2. Sistemas de coordenadas (Longitud del hueso)}
\textbf{Procedimiento y resultado:} Excelente resolución. Dibuja un gráfico representativo muy claro, calcula los catetos correctamente ($6$ y $8$) y aplica el Teorema de Pitágoras ($\sqrt{100}$) obteniendo de manera impecable la longitud de $10$ cm. \\
\textbf{Veredicto:} \textbf{Correcto} (2/2 pts).

\subsection*{3. Vectores (Fuerza resultante)}
\textbf{Procedimiento y resultado:} Entiende el concepto de suma vectorial y agrupa correctamente las componentes $\hat{i}$ y $\hat{j}$. Sin embargo, comete un \textbf{error de signo} evidente en la componente $\hat{j}$: en lugar de sumar $(25 + 35)\hat{j}$, escribe una resta $(25 - 35)\hat{j}$, lo que la lleva a un vector incorrecto de $30\hat{i} - 10\hat{j}$. \\
\textbf{Veredicto:} \textbf{Parcialmente correcto} (Planteamiento válido, error de signo al operar) (1/2 pts).

\subsection*{4. Potencias de diez (Notación científica y prefijo)}
\textbf{Procedimiento y resultado:} Expresa la medida decimal adecuadamente empleando una potencia de base diez ($12,5 \times 10^{-6}$). Sin embargo, la estudiante no completó la segunda parte de la instrucción, ya que omitió identificar e indicar explícitamente el prefijo del Sistema Internacional correspondiente (micrómetros, $\mu$m). \\
\textbf{Veredicto:} \textbf{Incompleto / Parcialmente correcto} (1/2 pts).

\subsection*{5. Sistemas de unidades (Conversión de densidad)}
\textbf{Procedimiento y resultado:} Su primer planteamiento utilizando una ecuación de factores de conversión formales presenta tachaduras y resulta muy confuso. Sin embargo, debajo de éste ejecuta un desglose aritmético perfecto: divide entre 1000 para transformar los gramos a kilogramos, y luego multiplica ese valor por $1.000.000$ para transformar el volumen. Su resultado final numérico ($1,030$) es totalmente exacto y correcto. \\
\textbf{Veredicto:} \textbf{Correcto} (2/2 pts).

\vspace{0.5cm}
\section*{Comentarios Generales y Sugerencias}
La estudiante demuestra una buena comprensión general de la lógica de conversión aritmética y la geometría. Se sugieren los siguientes aspectos a mejorar:
\begin{itemize}
    \item \textbf{Precaución con la ley de los signos:} Revisar cuidadosamente las operaciones algebraicas en sumas vectoriales para evitar restar cuando corresponde sumar (como ocurrió en la Pregunta 3).
    \item \textbf{Notación formal:} Se recomienda practicar la estructuración limpia de los factores de conversión en cadena para evitar confusiones en procedimientos futuros (Pregunta 5).
\end{itemize}

\end{document}